\writeSectionFrame{\presentationTitle}{Das Problem}
\frameWithImageOnly{\BILDERPFAD/problem1.jpg}
\note{
	Wir möchten Objekte klonen, also identische Kopien von Objekten anlegen. Wir möchten 
	also einen Prototypen erzeugen, non dem später Kopien erstellt werden können. Diese
	Kopien können wir dann noch mal leicht anpassen, so dass der Protoyp auch als Vorlage
	dienen kann. 
}

\begin{frame}{Kopieren von Objekten}
	\begin{center}
 		\includegraphics[scale=0.22]{\BILDERPFAD/problem2.png}
 	\end{center}
\end{frame}

\note{
 	Das Kopieren von Objekten von außen ist jedoch nicht ganz einfach, denn es gibt
 	Felder, die von außen nicht zugänglich sind für einen Kopierer. Nicht alle privaten
 	Felder sind über Getter zugänglich. 
}

\begin{frame}{Kopieren von Objekten}
	\begin{center}
 		\includegraphics[scale=0.25]{\BILDERPFAD/problem3.png}
 	\end{center}
\end{frame}

\note{
 	Es gibt
 	beim Kopieren noch ein Problem: Um eine Kopie anlegen zu können, muss ich die genaue 
 	Klasse eines Objekts kennen. Wenn wir nur mit Interfaces arbeiten, holen wir uns mit
 	dem Kopieren also eine weitere Abhängigkeit ins Boot.
}

\begin{frame}{Kopieren von Objekten}
	\begin{center}
 		\includegraphics[scale=0.25]{\BILDERPFAD/problem4.png}
 	\end{center}
\end{frame}

\note{
 	Also müssen die Objekte sich selbst kopieren.
}